\section{Reading the run laws off the ends}\label{sec:inverse}

The boundary layer is a renewal mass function, and a renewal mass function
determines its interarrival law. So the ends of the row identify the run laws.
For a fresh plus start one law is read off the atoms directly, since
$P_N(N)=T_\plus(N)$ by Proposition~\ref{prop:fresh-identity}. The other needs the
boundary layer at the plus end, and the remark after the proof shows why the
ratios at the minus end cannot replace the atoms.

\begin{proposition}[identification from the ends]\label{thm:inverse}
Assume (RV) for both signs.
\begin{enumerate}[(a)]
\item Let the start be fresh with sign $\plus$. Then $P_N(N)=T_\plus(N)$ for every
$N$, so the atoms determine $f_\plus$. For every fixed $j\ge1$,
\[
 \frac{P_N(N-j)}{P_N(N)}\to u_\minus(j),\qquad
 \frac{P_N(j)}{T_\minus(N)}\to u_\plus(j).
\]
The first limits determine $f_\minus$ through
\begin{equation}\label{eq:renewal-inverse}
 f_\minus(j)=u_\minus(j)-\sum_{i=1}^{j-1}f_\minus(i)\,u_\minus(j-i),\qquad j\ge1.
\end{equation}
In particular the family $(P_N)_{N\ge1}$ determines both run laws.
\item Let the preparation satisfy $Q_\plus(N)/T_\plus(N)\to q\in(0,\infty)$, and put
$\pi_\plus=\Prob(\Sigma=\plus)$. Then $\pi_\plus>0$, $\nu_\plus(n)=o(T_\plus(n))$,
and for every $j\ge1$
\[
 \frac{P_N(N-j)}{T_\plus(N)}\to\kappa_\plus(j)
 =\pi_\plus u_\minus(j)+\sum_{d=1}^{j}\nu_\minus(d)\,u_\minus(j-d).
\]
Hence $P_N(N-j)/P_N(N)\to\kappa_\plus(j)/q$, and when the preparation and
$q$ are known these limits determine $u_\minus$, hence $f_\minus$.
\end{enumerate}
\end{proposition}

The statement uses only the boundary layer at fixed $j$, so (LR) is not needed.
In (a) the second limit is a consistency check on the law of the plus runs and
is not needed for the identification. The informative point is the remark
below. Read through ratios alone, the minus end determines the law of the plus
runs only up to a geometric thinning, and the normalization $\sum_nf_\plus(n)=1$
does not remove this freedom.

\begin{proof}[Proof of Proposition~\ref{thm:inverse}(a)]
For a fresh plus start $P_N(N)=F^\plus_N(N)=T_\plus(N)$ by
Proposition~\ref{prop:fresh-identity}, and $T_\plus(N)>0$. The values $T_\plus(N)$ for
$N\ge1$ give $f_\plus(N)=T_\plus(N)-T_\plus(N+1)$. Theorem~\ref{thm:bigjump}(a)
with $\sgn_0=\plus$ and~\eqref{eq:kappa} give $P_N(N-j)/T_\plus(N)\to
\kappa^{(\plus)}_\plus(j)=u_\minus(j)$, which is the first limit. At the minus
end, Theorem~\ref{thm:bigjump}(a) with the signs exchanged gives
$P_N(j)/T_\minus(N)\to\kappa^{(\plus)}_\minus(j)$, and~\eqref{eq:kappa} with the
signs exchanged gives $\kappa^{(\plus)}_\minus(j)=u_\plus(j)$ for $j\ge1$. This
is the second limit.

For the recursion, let $j\ge1$ and split the sum defining $u_\minus(j)$
according to the length $i$ of the first run. The term $m=1$ is $f_\minus(j)$.
For $m\ge2$ the first run has length $1\le i\le j-1$ and the other $m-1$ runs
sum to $j-i\ge1$. Hence
\[
 u_\minus(j)=f_\minus(j)+\sum_{i=1}^{j-1}f_\minus(i)\sum_{m\ge1}f_\minus^{*m}(j-i)
 =f_\minus(j)+\sum_{i=1}^{j-1}f_\minus(i)\,u_\minus(j-i),
\]
which is~\eqref{eq:renewal-inverse}. It gives $f_\minus(1)=u_\minus(1)$, and
for $j\ge2$ it expresses $f_\minus(j)$ through $u_\minus(1),\dots,u_\minus(j)$ and
$f_\minus(1),\dots,f_\minus(j-1)$. By induction on $j$ the sequence $u_\minus$
determines $f_\minus$.
\end{proof}

\begin{remark}[what the minus end alone determines]\label{rem:thinning}
Let the start be fresh with sign $\plus$, and suppose only the ratios at the
two ends are read. At the minus end the second limit of
Proposition~\ref{thm:inverse}(a) gives $P_N(j)/P_N(i)\to u_\plus(j)/u_\plus(i)$, which
determines $(u_\plus(j))_{j\ge1}$ only up to a positive factor. The normalization
$\sum_nf_\plus(n)=1$ does not fix this factor. To see this, let $0<c\le1$, let $G$ be
independent of the runs with $\Prob(G=n)=c(1-c)^{n-1}$ for $n\ge1$, and let
$X^c_\plus$ be the sum of $G$ independent copies of $X_\plus$. With generating functions
$F(z)=\sum_nf_\plus(n)z^n$ and $U=1/(1-F)$, the law of $X^c_\plus$ has generating
function $F_c=cF/(1-(1-c)F)$, which equals one at $z=1$. Then
\[
 \frac1{1-F_c}=\frac{1-(1-c)F}{1-F}=1+c\,(U-1),
\]
so $u^c_\plus(j)=c\,u_\plus(j)$ for every $j\ge1$. In words, the renewals of
$X^c_\plus$ are those of $X_\plus$ kept independently with probability $c$. The
support of $X^c_\plus$ contains that of $X_\plus$, so (SPAN) is inherited, and
$\Prob(X^c_\plus\ge n)\sim T_\plus(n)/c$ by the tail asymptotics of random sums of
subexponential variables~\cite{ekm1997}, so (RV) holds with the same index.
Replacing $X_\plus$ by $X^c_\plus$ leaves $u_\minus$ and every ratio
$u_\plus(j)/u_\plus(i)$ unchanged. Hence the ratios at both ends do not
distinguish $f_\plus$ from the one-parameter family $f^c_\plus$, each member of
which sums to one. The atoms $P_N(N)=T_\plus(N)$ separate them, and so does the
normalization by $T_\minus(N)$ in Proposition~\ref{thm:inverse}(a), since
$P_N(j)/T_\minus(N)$ changes by the factor $c$. The bulk also separates them, since
the mean of $X^c_\plus$ is $\mu_\plus/c$ and $p$ changes.
\end{remark}

\begin{proof}[Proof of Proposition~\ref{thm:inverse}(b)]
Since $Q_\plus(N)\le\pi_\plus$ and $Q_\plus(N)>0$ for large $N$, we have $\pi_\plus>0$.
Next, $\nu_\plus(n)=Q_\plus(n)-Q_\plus(n+1)$, and by hypothesis
$Q_\plus(n)=qT_\plus(n)+o(T_\plus(n))$. Under (RV) for the plus sign
$T_\plus(n+1)/T_\plus(n)\to1$, so
$\nu_\plus(n)=q\bigl(T_\plus(n)-T_\plus(n+1)\bigr)+o(T_\plus(n))=o(T_\plus(n))$. Now
Theorem~\ref{thm:preparation}(d), which needs only (RV) for the plus sign, gives
$P_N(N-j)/T_\plus(N)\to\kappa_\plus(j)$ for every $j\ge1$, with $\kappa_\plus$ as
in~\eqref{eq:kappa-prep}.

Proposition~\ref{prop:first-run} gives $P_N(N)=Q_\plus(N)$, so
$P_N(N-j)/P_N(N)\to\kappa_\plus(j)/q$. Suppose the preparation and $q$ are
known. Then $\kappa_\plus(j)$ is known for every $j\ge1$, and since
$\nu_\minus(0)=0$ and $u_\minus(0)=1$,
\[
 \pi_\plus u_\minus(j)=\kappa_\plus(j)-\nu_\minus(j)-\sum_{d=1}^{j-1}\nu_\minus(d)\,u_\minus(j-d),
\]
which determines $u_\minus(j)$ by induction on $j$, because $\pi_\plus>0$.
Then~\eqref{eq:renewal-inverse} determines $f_\minus$.
\end{proof}

The hypothesis of (b) holds, for example, when $\nu_\plus=\pi_\plus f_\plus$ with
$\pi_\plus>0$ and $\nu_\minus$ arbitrary, since then $Q_\plus=\pi_\plus T_\plus$ and
$q=\pi_\plus$. The stationary preparation fails it, since its atom is of order
$NT_\plus(N)$ by Corollary~\ref{cor:stationary}(a).
